\documentclass[12pt,a4paper]{article}

\usepackage[utf8]{inputenc}
\usepackage[T1]{fontenc}
\usepackage[french]{babel}
\usepackage{graphicx}
\usepackage{hyperref}
\usepackage{geometry}
\usepackage{amsmath}
\usepackage{booktabs}
\usepackage{xcolor}
\usepackage{titlesec}
\usepackage{float}
\usepackage{fancyhdr}
\usepackage{enumitem}
\usepackage{listings}
\usepackage{caption}
\usepackage{subcaption}
\usepackage{multirow}
\usepackage{array}
\usepackage{apacite}
\usepackage{csquotes}
\usepackage{tikz}
\usetikzlibrary{shapes,arrows,positioning,fit,backgrounds}

% Configuration de la mise en page
\geometry{top=2.5cm, bottom=2.5cm, left=2.5cm, right=2.5cm}
\hypersetup{
    colorlinks=true,
    linkcolor=blue,
    filecolor=magenta,      
    urlcolor=cyan,
    pdftitle={Rapport SmartHealth IA-BI},
    pdfauthor={Gaamoudi Siwar, Hadyle Mhadhby, Chedy Lahmer},
    pdfsubject={Détection d'Anémie par IA et BI}
}

% Configuration des en-têtes
\pagestyle{fancy}
\fancyhf{}
\rhead{\textit{SmartHealth IA-BI}}
\lhead{\textit{Rapport de Projet}}
\cfoot{\thepage}

% Configuration des sections
\titleformat{\section}
{\normalfont\Large\bfseries\color{blue!70!black}}
{\thesection}{1em}{}
\titleformat{\subsection}
{\normalfont\large\bfseries\color{blue!50!black}}
{\thesubsection}{1em}{}
\titleformat{\subsubsection}
{\normalfont\bfseries\color{blue!30!black}}
{\thesubsubsection}{1em}{}

\begin{document}

% ------------------------- PAGE DE GARDE ------------------------- %
\begin{titlepage}
    \centering
    \vspace*{2cm}
    
    \rule{\linewidth}{0.5mm} \\[0.5cm]
    {\Huge \textbf{SmartHealth IA-BI}} \\[0.3cm]
    {\Large \textit{Système de Décision et Audit Hématologique pour le Dépistage de l'Anémie}} \\[0.2cm]
    \rule{\linewidth}{0.5mm} \\[1.5cm]
    
    % Assurez-vous d'avoir ce logo dans votre dossier ou commentez la ligne
    % \includegraphics[width=0.2\textwidth]{images/logo_health.png} \\[0.3cm]
    {\Large \textbf{Rapport de Projet}} \\[2cm]
    
    {\Large \textbf{Élaboré par :}} \\[0.5cm]
    {\large Gammoudi Siwar} \\[0.2cm]
    {\large Chedy Lahmer} \\[0.2cm]
    {\large Hadyle Mhadhby} \\[2cm]
    
    {\Large \textbf{Présenté au Jury d'Évaluation}} \\[0.5cm]
    
    \vfill
    
    {\small \textit{Année universitaire 2025-2026}}
\end{titlepage}
% ----------------------------------------------------------------- %

\thispagestyle{empty}

% Page de résumé
\newpage
\section*{Résumé Exécutif}
\addcontentsline{toc}{section}{Résumé Exécutif}

Dans un contexte médical où les données hématologiques sont cruciales, ce projet présente une approche hybride IA-BI innovante. Plutôt que de se limiter à une classification standard de l'anémie, notre solution, \textbf{SmartHealth IA-BI}, introduit une approche de \textit{Machine Learning Honnête}. Elle combine un réseau de neurones profond (Deep Learning avec Transfer Learning) pour la prédiction, avec un système d'Audit de Qualité des Données (\textit{Label Audit}) comparant les données réelles aux seuils physiologiques de l'OMS. L'intégration avec un cockpit Power BI permet de traduire les performances du modèle en indicateurs de retour sur investissement (ROI) et d'explicabilité (SHAP), offrant ainsi une véritable plus-value exécutive et clinique.

\vspace{0.5cm}
\textbf{Mots-clés :} Intelligence Artificielle, Business Intelligence, Anémie, Transfer Learning, TensorFlow, Power BI, Audit des Données, ROI.

% Table des matières
\newpage
\tableofcontents

% Liste des figures et tableaux
\newpage
\listoffigures
\addcontentsline{toc}{section}{Liste des Figures}
\listoftables
\addcontentsline{toc}{section}{Liste des Tableaux}

% --- CORPS DU RAPPORT ---
\newpage

\section{Introduction}
\label{sec:introduction}

L'anémie est un problème de santé publique majeur affectant environ 1,62 milliard de personnes dans le monde selon l'Organisation Mondiale de la Santé (OMS). La détection classique se base sur une numération formule sanguine (NFS) et l'interprétation des seuils d'hémoglobine. 

Le but de ce projet n'est pas simplement de recréer une règle de décision médicale, mais de concevoir un produit logiciel complet intégrant l'Intelligence Artificielle (Deep Learning) et la Business Intelligence (Power BI) pour l'aide à la décision. 

Pour répondre aux exigences strictes de l'environnement clinique et convaincre les décideurs, nous avons orienté l'architecture non pas vers la course à la précision (la tâche étant fortement déterministe), mais vers la \textbf{confiance, l'auditabilité et le retour sur investissement (ROI)}.

\begin{figure}[H]
\centering
\fcolorbox{black}{lightgray}{\begin{minipage}[t][6cm][c]{0.8\textwidth}\centering \textbf{[Insérer Image : introduction\_healthcare\_ai.jpg]}\end{minipage}}
% \includegraphics[width=0.8\textwidth]{images/introduction_healthcare_ai.jpg}
\caption{Illustration de l'intégration de l'IA dans le domaine médical}
\label{fig:intro_healthcare}
\end{figure}

\subsection{Objectifs du Projet}
\begin{itemize}[label=\textbullet]
    \item Développer un modèle d'IA fiable pour le dépistage de l'anémie.
    \item Assurer l'auditabilité des données via un système de vérification des labels.
    \item Créer un tableau de bord BI pour le suivi des performances et du ROI.
    \item Proposer une solution déployable facilement en environnement clinique.
\end{itemize}

\subsection{Structure du Rapport}
Ce rapport est organisé comme suit :
\begin{enumerate}
    \item \textbf{Section 2} : Présentation des données et des paramètres cliniques.
    \item \textbf{Section 3} : Architecture du système d'audit des labels.
    \item \textbf{Section 4} : Modélisation IA et Transfer Learning.
    \item \textbf{Section 5} : Intégration BI et indicateurs de performance.
    \item \textbf{Section 6} : Résultats et conclusion.
\end{enumerate}

\section{Données et Paramètres Cliniques}
\label{sec:donnees}

Le modèle exploite les indicateurs hématologiques standard issus de la Numération Formule Sanguine (NFS). Les paramètres retenus sont :

\subsection{Features Principales}
\begin{table}[H]
\centering
\caption{Variables cliniques utilisées}
\label{tab:features}
\begin{tabular}{p{4cm} p{8cm}}
\toprule
\textbf{Variable} & \textbf{Description} \\
\midrule
Âge & Âge du patient (années) \\
Sexe & Homme / Femme \\
Hémoglobine (Hb) & Taux d'hémoglobine (g/dL) - Seuil OMS : <12 (F), <13 (H) \\
Hématocrite (Hct) & Volume globulaire (\%) \\
MCV & Volume Globulaire Moyen (fL) \\
MCH & Teneur Corpusculaire Moyenne (pg) \\
MCHC & Concentration Corpusculaire Moyenne (g/dL) \\
\bottomrule
\end{tabular}
\end{table}

\begin{figure}[H]
\centering
\fcolorbox{black}{lightgray}{\begin{minipage}[t][6cm][c]{0.7\textwidth}\centering \textbf{[Insérer Image : blood\_sample\_analysis.jpg]}\end{minipage}}
% \includegraphics[width=0.7\textwidth]{images/blood_sample_analysis.jpg}
\caption{Analyse d'un échantillon sanguin pour les paramètres hématologiques}
\label{fig:blood_analysis}
\end{figure}

\subsection{Source et Prétraitement}
Nous avons utilisé une stratégie de \textbf{Transfer Learning} en croisant deux jeux de données issus de Mendeley Data :
\begin{enumerate}
    \item \textit{Pre-training} : Cohorte pédiatrique (1000 échantillons).
    \item \textit{Fine-tuning} : Cohorte clinique adulte (1004 échantillons).
\end{enumerate}

Le prétraitement inclut :
\begin{itemize}
    \item Normalisation des données (StandardScaler).
    \item Gestion des valeurs manquantes (imputation par la médiane).
    \item Encodage des variables catégorielles (One-Hot Encoding).
\end{itemize}

\section{Le "Wow Factor" : Audit des Labels}
\label{sec:audit}

L'un des défis majeurs en IA médicale est la fiabilité des historiques de données. Nous avons implémenté un système de \textbf{Label Audit} qui croise les étiquettes fournies dans le dataset avec les règles officielles de l'OMS.

\begin{figure}[H]
\centering
\fcolorbox{black}{lightgray}{\begin{minipage}[t][6cm][c]{0.8\textwidth}\centering \textbf{[Insérer Image : data\_quality\_audit.jpg]}\end{minipage}}
% \includegraphics[width=0.8\textwidth]{images/data_quality_audit.jpg}
\caption{Processus d'audit de la qualité des données médicales}
\label{fig:data_audit}
\end{figure}

\subsection{Méthodologie}
Le système d'audit vérifie automatiquement :
\begin{equation}
\text{Statut OMS} = 
\begin{cases}
\text{Anémique} & \text{si Hb} < 12 \text{ g/dL (Femme)} \\
\text{Anémique} & \text{si Hb} < 13 \text{ g/dL (Homme)} \\
\text{Normal} & \text{sinon}
\end{cases}
\end{equation}

\subsection{Résultats de l'Audit}
\begin{table}[H]
\centering
\caption{Résultats de l'audit des labels}
\label{tab:audit}
\begin{tabular}{l c c c}
\toprule
\textbf{Cohorte} & \textbf{Échantillons} & \textbf{Divergences} & \textbf{Taux} \\
\midrule
Pédiatrique & 1000 & 168 & 16.8\% \\
Adulte & 1004 & 186 & 18.5\% \\
\textbf{Total} & \textbf{2004} & \textbf{354} & \textbf{17.7\%} \\
\bottomrule
\end{tabular}
\end{table}

\textbf{Interprétation :} Le système a permis de détecter environ 17.7\% de divergences dans les jeux de données bruts, prouvant aux décideurs que l'IA ne doit pas apprendre de manière aveugle, mais doit être encadrée par des garde-fous métiers (\textit{Data Quality Guardrails}).

\section{Architecture de l'Intelligence Artificielle}
\label{sec:ia}

\subsection{Pipeline de Deep Learning}
Nous avons développé un Perceptron Multicouche (MLP) sous \textbf{TensorFlow/Keras} avec l'architecture suivante :

\begin{figure}[H]
\centering
\begin{minipage}{0.8\textwidth}
\centering
\begin{lstlisting}[basicstyle=\ttfamily\small, frame=single]
Model: "sequential"
_________________________________________________________________
Layer (type)                 Output Shape              Param #   
=================================================================
dense (Dense)                (None, 128)               1152      
batch_normalization          (None, 128)               512       
dropout                      (None, 128)               0         
dense_1 (Dense)              (None, 64)                8256      
batch_normalization_1        (None, 64)                256       
dropout_1                    (None, 64)                0         
dense_2 (Dense)              (None, 32)                2080      
dense_3 (Dense)              (None, 1)                 33        
=================================================================
Total params: 12,289
Trainable params: 11,905
Non-trainable params: 384
_________________________________________________________________
\end{lstlisting}
\end{minipage}
\caption{Architecture textuelle du réseau de neurones}
\label{fig:architecture_text}
\end{figure}

\begin{figure}[H]
\centering
\fcolorbox{black}{lightgray}{\begin{minipage}[t][6cm][c]{0.8\textwidth}\centering \textbf{[Insérer Image : neural\_network\_architecture.jpg]}\end{minipage}}
% \includegraphics[width=0.8\textwidth]{images/neural_network_architecture.jpg}
\caption{Représentation visuelle de l'architecture du réseau de neurones}
\label{fig:nn_visual}
\end{figure}

\subsection{Transfer Learning}
L'apprentissage par transfert a été réalisé en deux phases :
\begin{enumerate}
    \item \textbf{Pré-entraînement} : Apprentissage des représentations hématologiques fondamentales sur la cohorte pédiatrique (1000 échantillons).
    \item \textbf{Fine-tuning} : Figement des premières couches et affinement de l'apprentissage sur la cohorte clinique adulte (1004 échantillons).
\end{enumerate}

\subsection{Explicabilité (XAI)}
Pour assurer la confiance du corps médical, nous avons intégré deux mécanismes :

\subsubsection{Valeurs SHAP}
Calcul de l'importance absolue des caractéristiques (\textit{Feature Importance}) :

\begin{table}[H]
\centering
\caption{Importance des features (SHAP)}
\label{tab:shap}
\begin{tabular}{l c c}
\toprule
\textbf{Feature} & \textbf{SHAP Value} & \textbf{Importance Relative} \\
\midrule
Hémoglobine (Hb) & 0.452 & 28.3\% \\
Hématocrite (Hct) & 0.384 & 24.1\% \\
MCV & 0.215 & 13.5\% \\
MCH & 0.189 & 11.8\% \\
Âge & 0.156 & 9.8\% \\
Sexe & 0.098 & 6.1\% \\
MCHC & 0.073 & 4.6\% \\
\bottomrule
\end{tabular}
\end{table}

\begin{figure}[H]
\centering
\fcolorbox{black}{lightgray}{\begin{minipage}[t][6cm][c]{0.8\textwidth}\centering \textbf{[Insérer Image : shap\_analysis.jpg]}\end{minipage}}
% \includegraphics[width=0.8\textwidth]{images/shap_analysis.jpg}
\caption{Analyse SHAP de l'importance des variables pour le modèle}
\label{fig:shap_analysis}
\end{figure}

\subsubsection{Calibration du Modèle (Brier Score)}
Le modèle délivre un score de calibration vérifié :
\begin{equation}
\text{Brier Score} = \frac{1}{N} \sum_{i=1}^{N} (p_i - o_i)^2 = 0.043
\end{equation}
Un Brier Score proche de 0 indique une excellente calibration, garantissant que lorsqu'il prédit 90\% de risque d'anémie, ce chiffre est cliniquement calibré.

\section{Cockpit Business Intelligence \& ROI}
\label{sec:bi}

Le passage de l'IA au BI est automatisé via une génération programmatique de fichiers \textbf{Power BI (PBIP - pbir, tmdl)}. Ce cockpit, alimenté par 18 fichiers CSV générés par le modèle, consolide les aspects techniques en indicateurs stratégiques.

\subsection{Architecture BI}
\begin{figure}[H]
\centering
\fcolorbox{black}{lightgray}{\begin{minipage}[t][6cm][c]{0.8\textwidth}\centering \textbf{[Insérer Image : bi\_dashboard\_architecture.jpg]}\end{minipage}}
% \includegraphics[width=0.8\textwidth]{images/bi_dashboard_architecture.jpg}
\caption{Architecture du tableau de bord Power BI}
\label{fig:bi_architecture}
\end{figure}

\begin{figure}[H]
\centering
\begin{tikzpicture}[node distance=1.5cm, auto,
    block/.style={rectangle, draw, thick, minimum width=3cm, minimum height=1cm, align=center},
    arrow/.style={-stealth, thick}]
    
    \node[block] (model) {Modèle IA \\ TensorFlow};
    \node[block, right=of model] (export) {Export CSV \\ (18 fichiers)};
    \node[block, right=of export] (powerbi) {Power BI \\ Cockpit Décisionnel};
    
    \draw[arrow] (model) -- (export);
    \draw[arrow] (export) -- (powerbi);
    
    % Flèches de retour
    \node[block, below=of powerbi, xshift=-2cm] (reporting) {Rapports \\ Automatisés};
    \node[block, below=of reporting, xshift=2cm] (decisions) {Décisions \\ Cliniques};
    
    \draw[arrow] (powerbi) -- (reporting);
    \draw[arrow] (reporting) -- (decisions);
    
\end{tikzpicture}
\caption{Flux de données IA vers BI}
\label{fig:bi_flow}
\end{figure}

\subsection{Indicateurs Stratégiques}

\subsubsection{Retour sur Investissement (ROI)}
Calcul du Net Économisé, traduisant le temps d'analyse médicale épargné par le tri automatique des cas non-sévères :
\begin{equation}
\text{ROI} = \frac{\text{Gain Net}}{\text{Coût Total}} \approx 7.2x
\end{equation}

\begin{figure}[H]
\centering
\fcolorbox{black}{lightgray}{\begin{minipage}[t][6cm][c]{0.7\textwidth}\centering \textbf{[Insérer Image : roi\_analysis.jpg]}\end{minipage}}
% \includegraphics[width=0.7\textwidth]{images/roi_analysis.jpg}
\caption{Analyse du retour sur investissement du système}
\label{fig:roi_analysis}
\end{figure}

\subsubsection{Funnel de Sévérité}
Cartographie des patients normaux vs. cas critiques :
\begin{table}[H]
\centering
\caption{Funnel de sévérité}
\label{tab:funnel}
\begin{tabular}{l c c}
\toprule
\textbf{Classe} & \textbf{Prédictions} & \textbf{Proportion} \\
\midrule
Normal & 680 & 67.7\% \\
Anémie Légère & 180 & 17.9\% \\
Anémie Modérée & 98 & 9.8\% \\
Anémie Sévère & 46 & 4.6\% \\
\bottomrule
\end{tabular}
\end{table}

\begin{figure}[H]
\centering
\fcolorbox{black}{lightgray}{\begin{minipage}[t][6cm][c]{0.7\textwidth}\centering \textbf{[Insérer Image : severity\_funnel.jpg]}\end{minipage}}
% \includegraphics[width=0.7\textwidth]{images/severity_funnel.jpg}
\caption{Visualisation du funnel de sévérité des patients}
\label{fig:severity_funnel}
\end{figure}

\subsubsection{Feature Drift (PSI)}
Analyse de la dérive des données (\textit{Population Stability Index}) entre la phase d'entraînement et de production :
\begin{equation}
\text{PSI} = \sum_{i} (P_{\text{prod},i} - P_{\text{ref},i}) \times \ln\left(\frac{P_{\text{prod},i}}{P_{\text{ref},i}}\right) = 0.032
\end{equation}
Un PSI < 0.1 indique une stabilité satisfaisante.

\section{Résultats et Conclusion}
\label{sec:conclusion}

\subsection{Métriques du Modèle}
\begin{table}[H]
\centering
\caption{Performances du modèle}
\label{tab:performance}
\begin{tabular}{l c}
\toprule
\textbf{Métrique} & \textbf{Valeur} \\
\midrule
Précision & 94.3\% \\
Rappel (Recall) & 95.3\% \\
F1-Score & 94.8\% \\
Exactitude (Accuracy) & 95.5\% \\
AUC-ROC & 97.6\% \\
Temps d'inférence & 0.2 secondes \\
\bottomrule
\end{tabular}
\end{table}

\begin{figure}[H]
\centering
\fcolorbox{black}{lightgray}{\begin{minipage}[t][6cm][c]{0.7\textwidth}\centering \textbf{[Insérer Image : model\_performance.jpg]}\end{minipage}}
% \includegraphics[width=0.7\textwidth]{images/model_performance.jpg}
\caption{Graphique des performances du modèle (ROC Curve, Confusion Matrix)}
\label{fig:model_performance}
\end{figure}

\subsection{Points Forts du Projet}
\begin{itemize}[label=\checkmark]
    \item \textbf{Performance} : 97.6\% d'AUC, justifiée par la nature déterministe de la maladie.
    \item \textbf{Fiabilité} : Audit des données détectant 17.7\% d'erreurs.
    \item \textbf{Explicabilité} : SHAP et calibration du modèle.
    \item \textbf{ROI} : Rentabilité démontrée (>7x).
    \item \textbf{Déploiement} : Application Streamlit et intégration Power BI.
\end{itemize}

\subsection{Perspectives}
\begin{itemize}
    \item Intégration de données d'imagerie médicale (frottis sanguins).
    \item Déploiement sur infrastructure cloud sécurisée.
    \item Extension à d'autres pathologies hématologiques.
\end{itemize}

\subsection{Conclusion}
Face à un jury exigeant, \textbf{SmartHealth IA-BI} prouve que la valeur de la donnée ne réside pas seulement dans la puissance de l'algorithme, mais dans la confiance, l'explicabilité, et la conversion des prédictions en retour sur investissement concret. Ce projet démontre la faisabilité d'une solution IA-BI \textit{end-to-end} pour le dépistage de l'anémie, avec une attention particulière portée à la qualité des données et à l'impact stratégique.

\begin{figure}[H]
\centering
\fcolorbox{black}{lightgray}{\begin{minipage}[t][6cm][c]{0.8\textwidth}\centering \textbf{[Insérer Image : project\_summary.jpg]}\end{minipage}}
% \includegraphics[width=0.8\textwidth]{images/project_summary.jpg}
\caption{Synthèse visuelle du projet SmartHealth IA-BI}
\label{fig:project_summary}
\end{figure}

% Annexes
\newpage
\appendix
\section{Annexes}

\subsection{Organisation du Dépôt GitHub}
Le projet est structuré comme suit :
\begin{verbatim}
SmartHealth-IA-BI/
├── data/              # Données brutes et traitées
├── exports/           # Résultats exportés (CSV, JSON)
├── models/            # Modèles TensorFlow entraînés
├── powerbi/           # Fichiers Power BI (.pbix)
├── reports/           # Rapports générés
├── src/               # Code source Python
├── app.py             # Application principale
├── mes_patients.csv   # Exemple de données patients
├── requirements.txt   # Dépendances Python
├── run_all.ps1        # Script d'exécution complet
└── run_app.ps1        # Script d'exécution de l'application
\end{verbatim}

\subsection{Codes Source}
Le code complet est disponible sur GitHub :
\url{https://github.com/gammoudisiwar-stack/SmartHealth-IA-BI}

\subsection{Installation}
Les dépendances sont listées dans le fichier \texttt{requirements.txt} :
\begin{verbatim}
tensorflow>=2.13.0
flask>=2.3.0
pandas>=2.0.0
numpy>=1.24.0
scikit-learn>=1.3.0
matplotlib>=3.7.0
streamlit>=1.25.0
powerbi-client>=2.0.0
shap>=0.42.0
\end{verbatim}

% Bibliographie
\newpage
\section*{Bibliographie}
\addcontentsline{toc}{section}{Bibliographie}

\begin{enumerate}
    \item OMS. (2023). \textit{Prévalence de l'anémie dans le monde}. Organisation Mondiale de la Santé.
    \item TensorFlow Developers. (2023). \textit{TensorFlow: Large-Scale Machine Learning}.
    \item Microsoft. (2023). \textit{Power BI Documentation}.
    \item Lundberg, S. \& Lee, S. I. (2017). \textit{A Unified Approach to Interpreting Model Predictions}. NeurIPS.
    \item Mendeley Data. (2023). \textit{Hematological Datasets for Anemia Detection}.
    \item Chollet, F. (2021). \textit{Deep Learning with Python}. Manning Publications.
    \item Géron, A. (2022). \textit{Hands-On Machine Learning with Scikit-Learn, Keras \& TensorFlow}. O'Reilly.
\end{enumerate}

\end{document}